\chapter*{}
\noindent{\textbf{RESUMO}}

\noindent{Este trabalho desenvolveu e avaliou o GUY, uma extensão para o VS Code com foco inicial na linguagem Python, capaz de gerar o Grafo de Fluxo de Controle (GFC) e calcular a Complexidade Ciclomática para os exemplos analisados. Foram realizadas a análise de ferramentas correlatas, o levantamento de requisitos, a prototipagem, a implementação da extensão e a validação técnica com três exemplos em Python. Nos experimentos, a ferramenta gerou, respectivamente, \(V(G)=4\), com quatro caminhos independentes; \(V(G)=2\), com dois caminhos independentes; e \(V(G)=3\), com três caminhos independentes. Os resultados indicaram coerência entre os grafos produzidos, as estruturas de controle dos códigos analisados e a fórmula da métrica. Dentro desse escopo experimental, o GUY integrou a visualização interativa do GFC, o cálculo de \(V(G)\) e a navegação entre o grafo e o código-fonte no ambiente local do VS Code. Usos profissionais ou educacionais podem ser investigados em avaliações futuras, mas não foram medidos neste trabalho.
}

\vspace{\onelineskip}

\noindent{\textbf{Palavras-chave}: grafo de fluxo de controle; complexidade ciclomática; teste estrutural; VS Code.}

\vspace{\onelineskip}
